\section{Finite height laws and their temperature response}\label{sec:finite}

The temperature dependence needed in the proof is already visible on a
finite graph.  We first recall the comparison inequalities for the dual
heights, and then derive a lower bound for the derivative of a height
variance.  The lower bound will later distinguish a change of temperature
from a displacement along the critical trajectory.

\subsection{Conventions and finite comparisons}

For a finite primary height graph $G=(V,E)$, orient its edges once and
write $\nabla_e h=h_{e^+}-h_{e^-}$.  The height spacing is $2\pi$ and
its edge weight is
\begin{equation}\label{eq:bessel-weight}
 p_b(k)=e^{-b}I_k(b),\qquad
 \sum_{k\in\Z}p_b(k)e^{tk}=\exp\{b(\cosh t-1)\}.
\end{equation}
In particular, $p_b$ is the distribution of the difference of two
independent Poisson variables of mean $b/2$.  Each unoriented edge
contributes one factor.  The primary Laplacian $A_G$ is normalized by
\begin{equation}\label{eq:primary-energy}
 \langle v,A_Gv\rangle=\sum_{e\in E}(\nabla_ev)^2,
\end{equation}
where all discrete pairings use counting measure.  We never replace
$A_G$ by $A_G/4$.

A free connected component of the height law is taken modulo its common
$2\pi\Z$ translations.  A representative can be pinned at one vertex
when evaluating an invariant test.  If Gaussian observations confine the
constant, all its translates are instead summed.  On a torus we use
single-valued periodic heights, hence zero winding, with the same common
translation convention.  Its height space is
$2\pi\Z^{V}/(2\pi\Z\one)$; choosing mean-zero representatives means
projecting this lattice onto the mean-zero space, rather than requiring
each projected coordinate to remain an integer multiple of $2\pi$.  Additional primary pins always have value zero
unless an affine law is specified.

Fourier expansion of the spin interaction identifies the free spin
partition function with the divergence-free integer-current sum.  Rotating
currents onto the dual edges gives the height law with exterior face
height zero.  A spin insertion at two vertices changes the divergence to
two opposite unit sources; subtracting a unit current along a path between
them gives a height law with an integral connection on the crossed dual
edges.  This is the free-boundary duality of \cite[Lemma~2.1]{X}.
It applies to the edge convention in \eqref{eq:bessel-weight}; the factors
$e^b$ cancel in every partition ratio.

We shall use the following finite comparison rules.  A centered law may
include nonnegative quadratic penalties in real linear observations,
independent Gaussian observation noises, and centered linear equalities.
An affine law has the same precision and parallel constraint space but
may have an integral edge connection, nonzero equality values, or shifted
observation centers.

\begin{lemma}[Finite comparisons]\label{lem:comparisons}
For the finite height laws just described, increasing centered precision
decreases centered covariances and real moment generating functions.
The covariance in an affine law, also after a real linear tilt, dominates
the corresponding centered covariance.  A shifted-to-centered partition
ratio is at most one and decreases when the same positive precision is
added to both laws.  The latter monotonicity also holds with a fixed real
linear statistic inserted in the exponent of the affine numerator.

Centered invariant characteristic coefficients are positive and increase
with precision.  The Hessian of their logarithm at any phase is bounded
below by its Hessian at zero.  A characteristic test on a freely translated
component is invariant when the sum of its coefficients is an integer;
a real linear test on such a component must have coefficient sum zero.
For a centered real linear test $X$,
\begin{equation}\label{eq:laplace-lower}
                    \log\E e^X\ge\tfrac12\Var(X).
\end{equation}
The same centered moment and characteristic comparisons imply that, as
$b$ increases, centered covariances and real transforms increase, whereas
centered invariant characteristic coefficients decrease.
\end{lemma}

\begin{proof}
The lattice-Gaussian inequality underlying these comparisons is due
to Regev and Stephens-Davidowitz \cite{RSD}; related extensions to
annealed integer-valued Gaussian fields appear in
\cite[Appendix~D]{AHPS}.  The statements at fixed $b$ are \cite[Lemma~2.2]{X}.  We recall the
approximation which also gives the last assertion.  Replace each primary
edge by $N$ bonds in series, with normalized lattice-Gaussian bond weight
$(b/(2N))^{k^2}$.  On compact sets of real $t$ its transform is
\[
 \frac{\sum_k(b/(2N))^{k^2}e^{tk}}
      {\sum_k(b/(2N))^{k^2}}
 =1+\frac bN(\cosh t-1)+O(N^{-2}).
\]
The full chain therefore converges to \eqref{eq:bessel-weight}, including
all fixed exponential moments.  Matching endpoints, cycle constraints,
period constraints, and the additional observations are positive
quadratic forms or limits of such forms on the enlarged lattice.  The
finite lattice-Gaussian comparison inequalities apply there.  Increasing
$b$ decreases the bond precision $2\log(2N/b)$, proving the required
monotonicities before passage to the limit.  A spanning-forest
parametrization and full-chain exponential moments dominate the sums on
each fixed primary graph.  This justifies the limit, including fixed
polynomial insertions and real tilts.  Confined constants are dominated
by their observation penalties; unconfined constants are factored out.
\end{proof}

If $X=\sum_e w_e\nabla_eh$, its \emph{reference cost} is
$\sum_e w_e^2$, minimized over flows representing the same test.  For
$X=\langle h,f\rangle$ with $\sum_xf_x=0$, the minimum is
\begin{equation}\label{eq:reference-cost}
                 G_G(f)=\langle f,A_G^+f\rangle.
\end{equation}
The analogous definition uses a grounded Laplacian when primary vertices
are pinned.

\begin{lemma}[Uniform reference bounds]\label{lem:reference-cost}
Fix a compact interval $I\subset(0,\infty)$.  In every centered finite
height law above, for $b\in I$ and each fixed $r\ge1$,
\begin{equation}\label{eq:reference-moments}
 \Var(X)\le C_I\sum_e w_e^2,\qquad
 \E|X|^r\le C_{I,r}\Bigl(\sum_e w_e^2\Bigr)^{r/2}.
\end{equation}
For every fixed $t$, the transform
$\E\exp(tX/(\sum_e w_e^2)^{1/2})$ is bounded uniformly in the graph.
More generally, if $\max_e|tw_e|\le K$, then
\begin{equation}\label{eq:bounded-flow-mgf}
                \log\E e^{tX}\le C_{I,K}t^2\sum_e w_e^2.
\end{equation}
Zero tests are interpreted separately.
\end{lemma}

\begin{proof}
Release the centered endpoint, curl, period, and observation constraints
in the full-chain construction of Lemma~\ref{lem:comparisons}.  The
comparison bounds the transform by a product of independent whole-chain
transforms.  Their limiting logarithm is
$b\sum_e(\cosh(2\pi tw_e)-1)$, which is bounded by the right side of
\eqref{eq:bounded-flow-mgf} when the arguments are bounded.  The same
bound holds uniformly in sufficiently fine subdivision.  For a flow
normalized by its Euclidean norm every coefficient has modulus at most
one.  Fixed exponential moments then give \eqref{eq:reference-moments}
and the normalized transform bound.  This is the argument of
\cite[Lemma~2.3]{X}; its constants are uniform on $I$.
\end{proof}

\subsection{Poisson insertions and a finite-volume continuity estimate}

Retain the two Poisson variables $U_e,V_e$ for every oriented edge in
\eqref{eq:bessel-weight}, and impose
$U_e-V_e=\nabla_eh/(2\pi)$.  The joint unnormalized weight is
\[
 \prod_{e\in E}e^{-b}\frac{(b/2)^{U_e+V_e}}{U_e!V_e!}
\]
with the same height constraints and observation penalties as before.
Let $N=\sum_e(U_e+V_e)$ and $M=|E|$.  For $\ell\ge0$, write
$(N)_\ell=N(N-1)\cdots(N-\ell+1)$, with $(N)_0=1$.

\begin{lemma}[Factorial insertions]\label{lem:factorial}
In a centered finite height law, including centered primary pins and
centered observation penalties,
\begin{equation}\label{eq:factorial-bound}
                         \E_b(N)_\ell\le(bM)^\ell.
\end{equation}
If the observation centers are prescribed real values, the unnormalized
factorial insertion is bounded by $(bM)^\ell$ times the corresponding
zero-observation centered sum.
\end{lemma}

\begin{proof}
Expand $(N)_\ell$ as the sum over ordered choices of $\ell$ Poisson
units, recording their edges and signs.  Removing those units leaves a
partition sum with an integral edge connection and the same primary
height variables, pins, and observation penalties.  In particular the
remaining Poisson differences obey the shifted curl and period
constraints determined by that connection.  Every removed unit
contributes $b/2$.  Lemma~\ref{lem:comparisons} bounds each shifted sum
by its centered zero-observation counterpart.  Summing the $(2M)^\ell$
choices proves the assertion.  Repeated choices of the same edge cause
no difficulty: falling factorials are exactly the factors canceled by
removing several units from its Poisson factorial.
\end{proof}

For a fixed finite graph with observations $\eta$, let $\mu_b^\eta$
denote its unnormalized joint weight, including the Poisson variables,
and let $Z_b^0$ denote its centered zero-observation mass.  The graph may
have centered hard pins; these are included in both objects.

\begin{lemma}[Finite-volume parameter continuity]\label{lem:finite-perturbation}
Fix $b_0>0$ and a compact set $K$ of complex observation vectors on a
fixed finite list of observation shapes.  There are constants $c,C>0$
such that, for $\eta\in K$ and $|b-b_0|M\le c$,
\begin{equation}\label{eq:finite-tv}
 \frac{\|\mu_b^\eta-\mu_{b_0}^\eta\|_{\rm TV}}{Z_{b_0}^0}
                   \le C|b-b_0|M.
\end{equation}
Here total variation for complex weights means the sum of their absolute
values.  The constants can be uniform over any prescribed finite family
of graph topologies and observation shapes.  In particular, each
centered normalizer is nonzero on a disk of radius $c/M$ about $b_0$.
\end{lemma}

\begin{proof}
The parameter likelihood in the augmented variables is exactly
\begin{equation}\label{eq:poisson-likelihood}
       \exp\{-(b-b_0)M\}(b/b_0)^N.
\end{equation}
Expand the second factor in falling factorials.  For real observations,
Lemma~\ref{lem:factorial} bounds the absolute sum of its positive-order
terms, divided by $Z_{b_0}^0$, by
\[
 \sum_{\ell\ge1}\frac{|b-b_0|^\ell}{\ell!b_0^\ell}(b_0M)^\ell
                  =e^{|b-b_0|M}-1.
\]
The scalar exponential in \eqref{eq:poisson-likelihood} has the same
bound.  For $\eta=\eta_1+i\eta_2$, the absolute observation penalty is
its real-$\eta_1$ penalty times $e^{\|\eta_2\|^2/2}$, uniformly bounded
on $K$.  This proves \eqref{eq:finite-tv}.  At zero observation it gives
$|Z_b^0/Z_{b_0}^0-1|<1/2$ after decreasing $c$, hence nonvanishing.
All sums are analytic on compact parameter sets by the same expansion.
\end{proof}

\subsection{A strict response bound}

The upper bound in Lemma~\ref{lem:reference-cost} alone does not imply
that changing $b$ moves the long-scale height coefficient.  The following
inequality supplies the needed strict response using only finite graphs.

\begin{lemma}[Temperature response]\label{lem:response}
Let $G$ be a finite connected graph with the centered Bessel height law,
and let $f\ne0$ have coefficient sum zero.  Set
$X=\langle h,f\rangle$ and $G(f)=\langle f,A_G^+f\rangle$.
For $b$ in a compact interval $I\subset(0,\infty)$,
\begin{equation}\label{eq:temperature-response}
 \frac{\Var_b(X)^2}{4\pi^2b^2G(f)}
       \le \partial_b\Var_b(X)
       \le C_I|E|G(f).
\end{equation}
In particular, these inequalities hold on every zero-winding periodic
torus, with constants independent of its size.
\end{lemma}

\begin{proof}
Let $C_b$ be the covariance operator on mean-zero tests.  For each
oriented edge $e$, removing one positive Poisson unit gives the shifted
height weight
\[
 \prod_{d\in E}p_b\bigl(\nabla_dh/(2\pi)-\one_{\{d=e\}}\bigr).
\]
Let $r_e$ be its partition sum divided by the centered sum, and $v_e$
its mean of $X$.  Reversing $h$ identifies the negative-removal law with
its reflection, so its partition ratio is $r_e$ and its mean is $-v_e$.
By Lemma~\ref{lem:comparisons}, $0<r_e\le1$ and each shifted variance
is at least $\Var_b(X)$.

Differentiating the normalized augmented law inserts the score $N/b$;
the deterministic term $-|E|$ cancels.  Since the centered mean of $X$
is zero,
\[
 \partial_b\Var_b(X)=b^{-1}\Cov_b(X^2,N).
\]
Removing a positive or negative unit separately now gives the exact
identity
\begin{equation}\label{eq:response-insertions}
 \partial_b\Var_b(X)
   =\sum_{e\in E}r_e\bigl(\Var_{e,+}(X)-\Var_b(X)+v_e^2\bigr)
   \ge\sum_{e\in E}r_ev_e^2.
\end{equation}
The same removal with the insertion $X$ gives
\[
 \E_b[(U_e-V_e)X]=br_ev_e.
\]
As $\nabla_eh=2\pi(U_e-V_e)$, this is
\begin{equation}\label{eq:response-gradient}
                    \nabla_e(C_bf)=2\pi br_ev_e.
\end{equation}
Discrete integration by parts and Cauchy--Schwarz yield
\[
 \begin{split}
 \Var_b(X)^2
 &=\langle\nabla A_G^+f,\nabla C_bf\rangle^2\\
 &\le G(f)\sum_e|\nabla_eC_bf|^2
 \le4\pi^2b^2G(f)\sum_er_ev_e^2.
 \end{split}
\]
Together with \eqref{eq:response-insertions}, this proves the lower
bound.  For the upper bound, Lemma~\ref{lem:factorial} gives
$\E N^2\le(bM)^2+bM$, while
Lemma~\ref{lem:reference-cost} gives $\E X^4\le C_IG(f)^2$.
Cauchy--Schwarz in the score identity proves
$|\partial_b\Var_b(X)|\le C_IMG(f)$, since $M\ge1$ and $b\in I$.
The finite-volume analyticity established above justifies every
differentiation.  No bound on the inverse precision of a subdivided
Gaussian chain is used.
\end{proof}
